% ============================================================
% 纯答案册·课时15-2.7.1抛物线方程 —— 工具/答案抽册器.py 抽册生成件（勿手改；第三档＝纯答案单独成册）
% 源件（只读）：C:/提示词/工作区/M3-第2章量产0913/成卷/导学件/课时15-2.7.1抛物线方程/main.tex（md5 a32178f84d12372fb7461c6700148868）
%% 口径：题面不重印；逐题＝题号(印面号同正册)＋[答案]＋[详解]，值/详解照源件逐字
%       （钉值门硬断言：册值≡正册件内值≡值台账值tex，逐字全等）。册序＝正册装配序＝印面号 1..36 升序（同序同号）；键序断言基准＝题面库冻结 manifest canonical 键序。
%% 三档导出：整体 main-true／纯题 main-false（在产）｜本册＝纯答案册（本件）
%% 双档：ansbook-true.tex＝详解印本；ansbook-false.tex＝纯值速查本（详解不印）
% 编译：xelatex <壳> ×2，cwd＝本目录；sty＝qp-m3.sty 本地挂载（md5 7c3930362be8a0a2bdf21bbf8ac16573，锁校验）
% ============================================================
\documentclass[fontset=none]{ctexart}
\usepackage{qp-m3}
% —— 印面逐字符号路由补钉（承源件；− 走 TNR、▱ NSC 子块；禁改 sty 保 md5） ——
\xeCJKDeclareCharClass{Default}{"2212}%
\setCJKfamilyfont{nscblk}[Path=C:/提示词/工作区/字替对照-0909/variantF/fonts/,BoldFont=NSC-w600.ttf]{NSC-w500.ttf}%
\xeCJKDeclareSubCJKBlock{nscblk}{"25B1}%
\newcommand{\pxparallelogram}{{\CJKfamily{nscblk}▱}}%
% —— 断行微调（承母版实测档，三零门承重；\emergencystretch 不另设） ——
\xeCJKsetup{CJKglue={\hskip 0.10em plus 0.08em minus 0.05em}}%
% —— 纯答案册全线括线（长详解可跨栏断，承练习件全线括线口径；灰底盒不可跨栏断） ——
\ansblockgrayfalse
% —— 双档开关（false 档＝纯值速查：答案印、详解不印；件面开关，不动 sty 本体） ——
\newif\ifansbookdetail
\ifdefined\ansbookdetail
  \ifnum\ansbookdetail=0 \ansbookdetailfalse\else\ansbookdetailtrue\fi
\else
  \ansbookdetailtrue
\fi
% —— 页脚件名（承源件章名；件名＝<件型>答案册） ——
\renewcommand{\qpzhangming}{第二章\quad 平面解析几何}%
\renewcommand{\qpjianming}{导学件答案册}%

\begin{document}
\zhangtitle{第二章\quad 平面解析几何}
\jietitle{2.7.1 抛物线的标准方程}
\par\glueguard{1}\addvspace{2pt}\noindent\biaoqian{【纯答案册】}{\fangsong 题面不重印\quad 印面号与正册同号同序}\par\addvspace{2pt}

\begin{multicols}{2}
\emergencystretch=1em

\begin{ansblock}[2章-练-课时15-01]
% ans:2章-练-课时15-01
\ansitem{1}{不是。当定点 F 在直线 l 上时，动点轨迹是与 l 垂直的直线（过 F）。}
\ifansbookdetail
\ansnote{详解}{取\(l\)为\(x\)轴、\(F\)为原点．动点\(P(x,y)\)到\(F\)的距离为\(\sqrt{x^{2}+y^{2}}\)，到\(l\)的距离为\(|y|\)．\par\noindent “到\(F\)与到\(l\)距离相等”：\(\sqrt{x^{2}+y^{2}}=|y|\)，平方得\(x^{2}+y^{2}=y^{2}\)，即\(x=0\)．轨迹为整条\(y\)轴：过\(F\)且垂直于\(l\)的直线，不是抛物线（\(p=0\)时抛物线无意义）．验算：点\((0,5)\)到\(F\)距离5、到\(x\)轴距离5，相等合．}
\fi
\end{ansblock}

\begin{ansblock}[2章-练-课时15-03]
% ans:2章-练-课时15-03
\ansitem{2}{x²=−12y}
\ifansbookdetail
\ansnote{详解}{\(F(0,-3)\)不在直线\(y=3\)上，由定义轨迹为抛物线：焦点\(F(0,-3)\)，准线\(y=3\)．开口向下，设\(x^{2}=-2py\)（\(p>0\)）：\(p/2=3\)，\(p=6\)，方程\(x^{2}=-12y\)．验算：顶点\((0,0)\)到\(F\)距3、到准线距3合．}
\fi
\end{ansblock}

\begin{ansblock}[2章-练-课时15-05]
% ans:2章-练-课时15-05
\ansitem{3}{C}
\ifansbookdetail
\ansnote{详解}{左端为\(P(x,y)\)到定点\(F(0,p/2)\)的距离；右端\(|y+p/2|\)为\(P\)到定直线\(l\)：\(y=-p/2\)的距离．\(F(0,p/2)\)不在\(l\)上（\(p>0\)），由抛物线定义，轨迹为抛物线，故选C．}
\fi
\end{ansblock}

\begin{ansblock}[2章-练-课时15-14]
% ans:2章-练-课时15-14
\ansitem{4}{y²=16x。}
\ifansbookdetail
\ansnote{详解}{设\(M(x,y)\)．“\(|MF|\)比到\(l\)的距离小2”即\(|x+6|-|MF|=2\)．当\(x\geqslant -6\)时：\(x+6-\sqrt{(x-4)^{2}+y^{2}}=2\)，即\(\sqrt{(x-4)^{2}+y^{2}}=x+4\)（隐含\(x\geqslant -4\)，自然满足），平方得\((x-4)^{2}+y^{2}=(x+4)^{2}\)，化简为\(y^{2}=16x\)（此时\(x\geqslant 0\geqslant -4\)，与所设一致）；\par\noindent 当\(x<-6\)时：\(-x-6-\sqrt{(x-4)^{2}+y^{2}}=2\)，即\(\sqrt{(x-4)^{2}+y^{2}}=-x-4\)，平方同样得\(y^{2}=16x\)，但该抛物线上\(x\geqslant 0\)，与\(x<-6\)矛盾，此段无轨迹．故轨迹方程\(y^{2}=16x\)（顶点在原点、开口向右的抛物线，焦点\(F(4,0)\)、准线\(x=-4\)）．验算：\(y^{2}=16x\)上任一点\(|MF|=x+4\)，到直线\(x=-6\)的距离为\(x+6\)，差恰为2合．\par\noindent 辨析：本题须防“大2”误读——若误作“\(|MF|\)比到直线\(x=-6\)的距离大2”，则\(|MF|=|x+8|\)，平方化简为\(y^{2}=24(x+2)\)，为平移型抛物线（顶点\((-2,0)\)），与标准形不符；按原题“小2”，平移准线得\(|MF|\)等于到直线\(x=-4\)的距离，才落到标准抛物线\(y^{2}=16x\)．}
\fi
\end{ansblock}

\begin{ansblock}[2章-拓-课时15-15]
% ans:2章-拓-课时15-15
\ansitem{5}{D}
\ifansbookdetail
\ansnote{详解}{由定义，抛物线上点到焦点的距离等于它到准线\(x=-p/2\)的距离，其最小值在顶点处取得，为\(p/2\)．恒大于\(1\Leftrightarrow p/2>1\Leftrightarrow p>2\)，故选D．验算：\(p=2\)时顶点到焦点距离\(=1\)，不满足“恒大于”，开区间正确合．}
\fi
\end{ansblock}

\begin{ansblock}[2章-练-课时15-02]
% ans:2章-练-课时15-02
\ansitem{6}{y²=3x 或 y²=−3x 或 x²=3y 或 x²=−3y。}
\ifansbookdetail
\ansnote{详解}{焦准距即\(p\)：\(p=3/2\)，\(2p=3\)．开口未指定，四种标准方程并立：\(y^{2}=3x\)、\(y^{2}=-3x\)、\(x^{2}=3y\)、\(x^{2}=-3y\)．}
\fi
\end{ansblock}

\begin{ansblock}[2章-练-课时15-07]
% ans:2章-练-课时15-07
\ansitem{7}{(1) y²=8x；(2) y²=6x。}
\ifansbookdetail
\ansnote{详解}{(1)焦点\((2,0)\)在\(x\)正半轴：\(p/2=2\)，\(p=4\)，\(2p=8\)，方程\(y^{2}=8x\)．(2)准线\(x=-3/2\)得焦点在\(x\)正半轴：\(p/2=3/2\)，\(p=3\)，\(2p=6\)，方程\(y^{2}=6x\)．验算：(1)准线\(x=-2\)合；(2)焦点\((3/2,0)\)合．}
\fi
\end{ansblock}

\begin{ansblock}[2章-拓-课时15-07]
% ans:2章-拓-课时15-07
\ansitem{8}{(a,0)（a≠0）。}
\ifansbookdetail
\ansnote{详解}{按标准形\(y^{2}=2px\)对照：\(2p=4a\)，\(p=2a\)，焦点\((p/2,0)=(a,0)\)．\(a>0\)时开口向右，焦点\((a,0)\)；\(a<0\)时开口向左，焦点仍为\((a,0)\)（即\(x\)轴负半轴上）．故焦点\((a,0)\)（\(a\neq 0\)）．验算：\(a=-1\)时\(y^{2}=-4x\)，焦点\((-1,0)\)合．}
\fi
\end{ansblock}

\begin{ansblock}[2章-拓-课时15-08]
% ans:2章-拓-课时15-08
\ansitem{9}{(1) 焦点 (0,1/8)，准线 y=−1/8；(2) 焦点 (0,1/(4a))，准线 y=−1/(4a)（a≠0）。}
\ifansbookdetail
\ansnote{详解}{(1)\(y=2x^{2}\)即\(x^{2}=(1/2)y=2\cdot (1/4)y\)：\(2p=1/2\)，\(p=1/4\)，开口向上，焦点\((0,1/8)\)，准线\(y=-1/8\)．\par\noindent (2)\(y=ax^{2}\)即\(x^{2}=(1/a)y=2\cdot [1/(2a)]y\)：\(2p=1/a\)，\(p=1/(2a)\)．\(a>0\)时开口向上，焦点\((0,1/(4a))\)，准线\(y=-1/(4a)\)；\(a<0\)时开口向下，焦点、准线表达式不变（\(1/(4a)<0\)，准线在\(x\)轴上方）．验算：\(a=2\)时与(1)一致，\(1/(4\times 2)=1/8\)合．}
\fi
\end{ansblock}

\begin{ansblock}[2章-拓-课时15-10]
% ans:2章-拓-课时15-10
\ansitem{10}{(0,−1/3)；y=1/3。}
\ifansbookdetail
\ansnote{详解}{\(3x^{2}+4y=0\)即\(x^{2}=-(4/3)y=-2\cdot (2/3)y\)：\(2p=4/3\)，\(p=2/3\)，开口向下，焦点\((0,-p/2)=(0,-1/3)\)，准线\(y=p/2=1/3\)．验算：顶点到焦点、到准线等距\(1/3\)合．}
\fi
\end{ansblock}

\begin{ansblock}[2章-拓-课时15-11]
% ans:2章-拓-课时15-11
\ansitem{11}{x²=8y}
\ifansbookdetail
\ansnote{详解}{准线\(y=-2\)在\(x\)轴下方，开口向上：设\(x^{2}=2py\)（\(p>0\)），\(-p/2=-2\)得\(p=4\)，\(2p=8\)，方程\(x^{2}=8y\)．验算：焦点\((0,4)\)、准线\(y=-2\)合．}
\fi
\end{ansblock}

\begin{ansblock}[2章-练-课时15-08]
% ans:2章-练-课时15-08
\ansitem{12}{y²=24x}
\ifansbookdetail
\ansnote{详解}{焦点在\(x\)正半轴，开口向右，设\(y^{2}=2px\)（\(p>0\)），准线\(x=-p/2\)．准线与\(y\)轴距离\(|-p/2|=6\)得\(p=12\)，\(2p=24\)，方程\(y^{2}=24x\)．验算：焦点\((6,0)\)、准线\(x=-6\)，与\(y\)轴距离均为6合．}
\fi
\end{ansblock}

\begin{ansblock}[2章-练-课时15-04]
% ans:2章-练-课时15-04
\ansitem{13}{y²=−16x}
\ifansbookdetail
\ansnote{详解}{焦点在\(x\)轴，设\(y^{2}=mx\)（\(m\neq 0\)）．代入\((-4,-8)\)：\(64=m\times (-4)\)，\(m=-16\)，方程\(y^{2}=-16x\)（开口向左，与横坐标为负相符）．}
\fi
\end{ansblock}

\begin{ansblock}[2章-练-课时15-09]
% ans:2章-练-课时15-09
\ansitem{14}{y²=8x}
\ifansbookdetail
\ansnote{详解}{设\(y^{2}=2px\)（\(p>0\)）．代入\((2,-4)\)：\(16=4p\)，\(p=4\)，\(2p=8\)，方程\(y^{2}=8x\)．验算：焦点\((2,0)\)、准线\(x=-2\)；\(M(2,-4)\)到焦点距离4、到准线距离4合．}
\fi
\end{ansblock}

\begin{ansblock}[2章-练-课时15-11]
% ans:2章-练-课时15-11
\ansitem{15}{y²=8x 或 x²=−y}
\ifansbookdetail
\ansnote{详解}{\(A(2,-4)\)在第四象限，开口向右或向下．开口向右：设\(y^{2}=2px\)（\(p>0\)），\((-4)^{2}=2p\times 2\)得\(p=4\)，方程\(y^{2}=8x\)；开口向下：设\(x^{2}=-2py\)（\(p>0\)），\(4=-2p\times (-4)\)得\(p=1/2\)，方程\(x^{2}=-y\)．综上\(y^{2}=8x\)或\(x^{2}=-y\)．验算：\(4^{2}=8\times 2\)合；\(2^{2}=-(-4)\)合．}
\fi
\end{ansblock}

\begin{ansblock}[2章-拓-课时15-12]
% ans:2章-拓-课时15-12
\ansitem{16}{D}
\ifansbookdetail
\ansnote{详解}{椭圆\(x^{2}/2+y^{2}=1\)：\(a^{2}=2\)、\(b^{2}=1\)，\(c^{2}=1\)，焦点\((\pm 1,0)\)，对称中心\(O(0,0)\)．以\(O\)为顶点、焦点为\((1,0)\)或\((-1,0)\)：\(p/2=1\)，\(2p=4\)，\(y^{2}=4x\)或\(y^{2}=-4x\)，故选D．验算：\(y^{2}=\pm 4x\)的焦点为\((\pm 1,0)\)合．}
\fi
\end{ansblock}

\begin{ansblock}[2章-拓-课时15-13]
% ans:2章-拓-课时15-13
\ansitem{17}{D}
\ifansbookdetail
\ansnote{详解}{由\((2\sqrt{2})^{2}=2px_0\)得\(x_0=4/p\)．\(M\)到准线的距离\(=x_0+p/2=4/p+p/2=3\)，得\(p^{2}-6p+8=0\)，\(p=2\)或\(p=4\)，抛物线\(y^{2}=4x\)或\(y^{2}=8x\)，故选D．验算：\(p=2\)时\(x_0=2\)，\(M(2,2\sqrt{2})\)到准线\(x=-1\)距离3合；\(p=4\)时\(x_0=1\)，\(M(1,2\sqrt{2})\)到准线\(x=-2\)距离3合．}
\fi
\end{ansblock}

\begin{ansblock}[2章-练-课时15-06]
% ans:2章-练-课时15-06
\ansitem{18}{3}
\ifansbookdetail
\ansnote{详解}{由抛物线定义，抛物线上点到焦点的距离与到准线的距离相等，故\(M\)到准线的距离为\(|MF|=3\)．}
\fi
\end{ansblock}

\begin{ansblock}[2章-练-课时15-10]
% ans:2章-练-课时15-10
\ansitem{19}{M(6,6√2) 或 M(6,−6√2)。}
\ifansbookdetail
\ansnote{详解}{\(y^{2}=12x\)：\(2p=12\)，\(p=6\)，准线\(x=-3\)．由定义\(|MF|=x_M+3=9\)，得\(x_M=6\)；代入得\(y^{2}=72\)，\(y=\pm 6\sqrt{2}\)．故\(M(6,6\sqrt{2})\)或\(M(6,-6\sqrt{2})\)（两点均在抛物线上，关于\(x\)轴对称）．验算：\(\sqrt{(6-3)^{2}+72}=\sqrt{81}=9\)合．}
\fi
\end{ansblock}

\begin{ansblock}[2章-练-课时15-12]
% ans:2章-练-课时15-12
\ansitem{20}{x=0；1/2。}
\ifansbookdetail
\ansnote{详解}{设\(P(x,y)\)：\(\sqrt{x^{2}+(y-1)^{2}}+|y+1|=3\)．①\(y>2\)或\(y<-4\)（\(|y+1|>3\)）：左边\(>3\)，无解；②\(-1\leqslant y\leqslant 2\)：\(\sqrt{x^{2}+(y-1)^{2}}=2-y\)，平方得\(x^{2}=-2y+3\)（\(-1\leqslant y\leqslant 3/2\)），对称轴\(x=0\)，此时\(|PF|=2-y\geqslant 1/2\)；\par\noindent ③\(-4\leqslant y<-1\)：\(\sqrt{x^{2}+(y-1)^{2}}=y+4\)，平方得\(x^{2}=10y+15\)（\(-3/2\leqslant y<-1\)），对称轴\(x=0\)，此时\(|PF|=y+4\geqslant 5/2\)．综上，对称轴方程\(x=0\)，\(|PF|_{\min}=1/2\)．验算：②段平方展开\(x^{2}+(y-1)^{2}=(2-y)^{2}\)得\(x^{2}=-2y+3\)合；②段与③段区间衔接经复核无缝合．}
\fi
\end{ansblock}

\begin{ansblock}[2章-拓-课时15-14]
% ans:2章-拓-课时15-14
\ansitem{21}{答案见解析：x 的系数为正数且越大时，抛物线的开口向右且开口越大。}
\ifansbookdetail
\ansnote{详解}{三条抛物线同以\(x\)轴为对称轴、顶点在原点、开口向右．作图观察（或取同一\(y\)值比较横坐标\(x=y^{2}/(2p)\)）：系数1、2、4递增时，同一高度处横坐标\(y^{2}/(2p)\)递减，曲线横向收拢更快、开口更大．故：\(x\)的系数为正且越大，开口向右且开口越大．}
\fi
\end{ansblock}

\begin{ansblock}[2章-练-课时15-13]
% ans:2章-练-课时15-13
\ansitem{22}{B}
\ifansbookdetail
\ansnote{详解}{以顶端\(O\)为原点建系，设\(x^{2}=-2py\)（\(p>0\)）．由对称性\(D(15,h)\)（\(h<0\)），\(B(30,h-150)\)．联立：\(225=-2ph\)，\(900=-2p(h-150)\)．两式相除：\(4=(h-150)/h\)，得\(h=-50\)；代入\(225=-2p\times (-50)\)得\(p=2.25\)．顶端\(O\)到\(AB\)的距离\(=|h|+150=200\)(m)，故选B．验算：\(p=2.25\)，\(2p=4.5\)，\(900=4.5\times 200\)合．}
\fi
\end{ansblock}

\begin{ansblock}[2章-练-课时15-16]
% ans:2章-练-课时15-16
\ansitem{23}{不能安全通过隧道。}
\ifansbookdetail
\ansnote{详解}{以抛物线顶点为原点、对称轴为\(y\)轴建系．由图：拱高3m、半宽3m，抛物线过\(A(3,-3)\)；矩形部分高2m、总高（顶点至地面）5m、地面\(y=-5\)．设\(x^{2}=-2py\)（\(p>0\)）：\(9=-2p\times (-3)\)，\(2p=3\)，方程\(x^{2}=-3y\)．\par\noindent 车居中：\(x=\pm 1.5\)处拱高\(y=-(1/3)\times 2.25=-0.75\)，净空\(=5-0.75=4.25\)(m)．\(4.25<4.5\)（车与集装箱总高），即使集装箱处于隧道正中，总高仍高于该处拱高，故不能安全通过．验算：\((\pm 1.5)^{2}=-3y\)得\(y=-0.75\)合；\(2p=3\)时准线\(y=3/4\)合．}
\fi
\end{ansblock}

\begin{ansblock}[2章-拓-课时15-01]
% ans:2章-拓-课时15-01
\ansitem{24}{桥拱所在抛物线 y=−(1/80)x²；溢流孔所在抛物线分别为 y=−(5/36)(x+14)²、y=−(5/36)(x+7)²、y=−(5/36)(x−7)²、y=−(5/36)(x−14)²；交点 A((140/13),−(245/169))、B(10,−5/4)、C((70/13),−(245/676))。}
\ifansbookdetail
\ansnote{详解}{设桥拱\(y=ax^{2}\)．由图注，桥拱过点\((20,-5)\)：\(-5=a\times 20^{2}\)，\(a=-1/80\)，桥拱\(y=-(1/80)x^{2}\)．①\par\noindent 设\(A\)所在溢流孔\(y=a_1(x-14)^{2}\)（顶点\((14,0)\)），亦过\((20,-5)\)：\(-5=a_1\times 36\)，\(a_1=-5/36\)．4孔轮廓相同、顶点\((\pm 7,0)\)、\((\pm 14,0)\)：\(y=-(5/36)(x\mp 7)^{2}\)、\(y=-(5/36)(x\mp 14)^{2}\)．②③\par\noindent 联立①②：\(400(x-14)^{2}=36x^{2}\)（两边乘2880），即\(91x^{2}-2800x+19600=0\)，\(\Delta=705600\)，\(\sqrt{\Delta}=840\)，\(x=20\)（\(y=-5\)）或\(x=140/13\)（\(y=-245/169\)），取交点\(A(140/13,-245/169)\)；联立①③：\(400(x-7)^{2}=36x^{2}\)，即\(91x^{2}-1400x+4900=0\)，\(\Delta=176400\)，\(\sqrt{\Delta}=420\)，\(x=10\)（\(y=-5/4\)）或\(x=70/13\)（\(y=-245/676\)），得\(C(70/13,-245/676)\)、\(B(10,-5/4)\)．验算：\(x=140/13\)：\(y=-(1/80)(140/13)^{2}=-19600/(80\times 169)=-245/169\)合；\(x=10\)：\(-100/80=-5/4\)合；\(x=70/13\)：\(-4900/(80\times 169)=-245/676\)合．}
\fi
\end{ansblock}

\begin{ansblock}[2章-拓-课时15-09]
% ans:2章-拓-课时15-09
\ansitem{25}{(1) y=−2x+6；(2) 6m。}
\ifansbookdetail
\ansnote{详解}{(1)\(|BF|=p/2=1/2\)，灯柱\(CD\)过焦点\(F(1/2,0)\)且\(A\)到灯柱距离1.5，得\(x_A=1/2+3/2=2\)；代入\(y^{2}=2x\)得\(y_A=2\)（取上半支），\(A(2,2)\)．设\(A\)处切线\(y-2=k(x-2)\)：与\(y^{2}=2x\)联立消\(x\)得\(ky^{2}-2y+4-4k=0\)，相切\(\Delta=4-4k(4-4k)=0\)，即\(16k^{2}-16k+4=0\)，\(k=1/2\)．灯罩轴线与切线垂直：斜率\(-2\)，方程\(y-2=-2(x-2)\)，即\(y=-2x+6\)．\par\noindent (2)路宽\(|DH|=10\)，灯罩轴线过路面中线，轴线与路面交点到\(y\)轴距离\(1/2+5=11/2\)．对\(y=-2x+6\)，\(x=11/2\)时\(y=-5\)，交点\((11/2,-5)\)，\(|FD|=5\)；又\(x=1/2\)代入\(y^{2}=2x\)得\(y=\pm 1\)，\(|CF|=1\)．灯柱高\(|CD|=|DF|+|FC|=6\)(m)．验算：\(16k^{2}-16k+4=4(2k-1)^{2}=0\)合；\(|CD|=6\)合．}
\fi
\end{ansblock}

\begin{ansblock}[2章-拓-课时15-03]
% ans:2章-拓-课时15-03
\ansitem{26}{BCD}
\ifansbookdetail
\ansnote{详解}{A：\(D_1D\perp\)底面，\(DN=\sqrt{MN^{2}-MD^{2}}=\sqrt{4-1}=\sqrt{3}\)，\(N\)的轨迹是以\(D\)为圆心、\(\sqrt{3}\)为半径的圆．设\(MN\)中点\(P\)、\(MD\)中点\(Q\)，则\(PQ\parallel ND\)且\(PQ=\sqrt{3}/2\)，\(P\)的轨迹是以\(Q\)为圆心\(\sqrt{3}/2\)为半径的圆，面积\(\pi\times 3/4=3\pi/4\neq \pi\)，A错；\par\noindent B：\(N\)到\(B_1B\)的距离\(=NB\)（\(B_1B\perp\)底面），即\(N\)到定点\(B\)与定直线\(DC\)距离相等（\(B\)不在\(DC\)上），轨迹为抛物线，B对；\par\noindent C：\(AB\parallel D_1C_1\)，\(D_1N\)与\(D_1C_1\)所成角\(\pi/3\)，即\(D_1N\)为以\(D_1C_1\)为轴、轴截面顶角\(2\pi/3\)的圆锥母线，底面\(ABCD\)与轴\(D_1C_1\)平行，交线为双曲线，C对；\par\noindent D：\(D_1N\)为以\(D_1B\)为轴、顶角\(\pi/3\)的圆锥母线，平面\(ABCD\)与轴\(D_1B\)斜交，交线为椭圆，D对．故选BCD．}
\fi
\end{ansblock}

\begin{ansblock}[2章-拓-课时15-04]
% ans:2章-拓-课时15-04
\ansitem{27}{BCD}
\ifansbookdetail
\ansnote{详解}{以\(D\)为原点、\(DA\)、\(DC\)、\(DD_1\)为\(x\)、\(y\)、\(z\)轴建系：\(A(2,0,0)\)，\(C(0,2,0)\)，\(B(2,2,0)\)，\(M(1,0,2)\)．\par\noindent A：\(\overrightarrow{AC}=(-2,2,0)\)，\(\overrightarrow{BM}=(-1,-2,2)\)，\(|\cos\langle\overrightarrow{AC},\overrightarrow{BM}\rangle|=|\overrightarrow{AC}\cdot\overrightarrow{BM}|/(|\overrightarrow{AC}||\overrightarrow{BM}|)=|2-4+0|/(2\sqrt{2}\times 3)=\sqrt{2}/6\neq\sqrt{2}/3\)，A错；\par\noindent B：\(S=(1/2)|AC_1|\cdot h=\sqrt{3}\)，\(|AC_1|=2\sqrt{3}\)，得\(h=1\)：\(P\)在以\(AC_1\)为轴、半径1的圆柱面上，该圆柱面与侧面\(BCC_1B_1\)（与轴斜交）的交线为椭圆的一部分，B对；\par\noindent C：\(P\)在侧面\(BCC_1B_1\)上时，到\(C_1D_1\)的距离\(=|PC_1|\)，条件即\(|PC_1|\)等于\(P\)到直线\(BC\)的距离——到定点与定直线等距且\(C_1\)不在\(BC\)上，轨迹为抛物线的一部分，C对；\par\noindent D：取\(AD\)中点\(H\)，\(MH=2\)且\(MH\perp\)底面；当交线与\(AC\)平行（\(HN\perp AC\)）时二面角最小，截面为过\(M\)、\(N\)及四棱中点的正六边形，边长\(\sqrt{2}\)，面积\(6\times(\sqrt{3}/4)\times 2=3\sqrt{3}\)，D对．故选BCD．}
\fi
\end{ansblock}

\begin{ansblock}[2章-拓-课时15-05]
% ans:2章-拓-课时15-05
\ansitem{28}{C}
\ifansbookdetail
\ansnote{详解}{折叠不变量：\(PB=AB\)，\(PD=AD\)，故\(PB+PD=10>BD=8\)，\(P\)在以\(B\)、\(D\)为焦点、长轴10、焦距8的椭圆上（不在长轴端点）．\par\noindent \(\triangle BPD\cong\triangle BCD\)（三边相等），得\(\angle PBD=\angle CBD\)．过\(P\)作\(PE\perp BD\)于\(E\)，连\(CE\)：\(\triangle BPE\cong\triangle BCE\)，得\(CE\perp BD\)且\(PE=CE\)，故\(BD\perp\)平面\(PCE\)．\(V_{P-BCD}=V_{B-PCE}+V_{D-PCE}=(1/3)S_{\triangle PCE}\cdot BD=(8/3)S_{\triangle PCE}\)．取\(PC\)中点\(F\)，\(EF\perp PC\)（\(PE=CE\)）：\(S_{\triangle PCE}=(1/2)\cdot PC\cdot EF=\sqrt{PE^{2}-1}\)（\(PC=2\)得\(CF=1\)）．\(PE\)的最大值为椭圆短半轴\(\sqrt{5^{2}-4^{2}}=3\)，故\(S_{\triangle PCE}\)最大\(=\sqrt{9-1}=2\sqrt{2}\)．\(V_{\max}=(8/3)\times 2\sqrt{2}=16\sqrt{2}/3\)，故选C．验算：\(EF=\sqrt{PE^{2}-CF^{2}}=\sqrt{PE^{2}-1}\)合；\(PE=3\)时\(EF=2\sqrt{2}\)，\(S=(1/2)\times 2\times 2\sqrt{2}=2\sqrt{2}\)合．}
\fi
\end{ansblock}

\begin{ansblock}[2章-拓-课时15-06]
% ans:2章-拓-课时15-06
\ansitem{29}{BCD}
\ifansbookdetail
\ansnote{详解}{过\(P\)作\(PO\perp\)平面\(ABCD\)于\(O\)，\(OH\perp AD\)于\(H\)．\(AD\perp\)平面\(POH\)，故\(AD\perp PH\)，\(\angle PHO=\alpha\)；又\(\angle PBO=\beta\)，\(\alpha=\beta\)得\(OH=OB\)．\par\noindent \(O\)的轨迹：底面内到定点\(B\)与定直线\(AD\)等距的点集，即以\(B\)为焦点、\(AD\)为准线的抛物线在四边形\(ABCD\)内（含边界）的部分；\(P\)的轨迹为其上平移（以\(B_1\)为焦点、\(A_1D_1\)为准线），是抛物线的一部分而非整条，A错；\par\noindent B：\(|PB|^{2}=(2\sqrt{2})^{2}+|OB|^{2}\)，\(O\)为\(AB\)中点时\(|OB|_{\min}=1\)，故\(|PB|_{\min}=\sqrt{8+1}=3\)，B对；\par\noindent C：\(PA_1\)与\(CD\)所成角即\(OA\)与\(AB\)所成角\(\angle OAB\)．以\(AB\)中点\(M\)为原点建系，抛物线\(y^{2}=4x\)，\(A(-1,0)\)．设切线\(x=ty-1\)：代入得\(y^{2}-4ty+4=0\)，\(\Delta=16t^{2}-16=0\)，\(t=\pm 1\)，取\(t=1\)（\(O\)在四边形内），\(\angle OAB=\pi/4\)，最大值\(\pi/4\)，C对；\par\noindent D：\(V_{P-A_1BC_1}=V_{B-PA_1C_1}=(1/3)S_{\triangle PA_1C_1}\cdot BB_1\)，\(A_1C_1=2\sqrt{2}\)．\(P\)到\(A_1C_1\)距离最大\(\Leftrightarrow O\)到\(AC\)距离最大：\(O\)为\(AB\)中点时最大，值\((1/4)|BD|=\sqrt{2}/2\)．\(V_{\max}=(2\sqrt{2}/3)\times (1/2)\times 2\sqrt{2}\times (\sqrt{2}/2)=2\sqrt{2}/3\)，D对．故选BCD．验算：\(y^{2}=4x\)（焦点\(B(1,0)\)，准线\(x=-1\)即\(AD\)所在直线）：\(|OB|\)与\(|OH|\)分别为到焦点、准线的距离，等距即轨迹合；切线\(y=x+1\)过\(A(-1,0)\)，斜率1，\(\angle OAB=45^{\circ}\)合．}
\fi
\end{ansblock}

\begin{ansblock}[2章-拓-课时15-02]
% ans:2章-拓-课时15-02
\ansitem{30}{A}
\ifansbookdetail
\ansnote{详解}{\(y^{2}=8x\)：\(p=4\)，准线\(x=-2\)，焦点\(F(2,0)\)．\(P\)在准线\(x=-2\)上，又在\(x-y+6=0\)上：\(P(-2,4)\)．\(k_{PF}=(0-4)/(2+2)=-1\)；由特征(3)\(PF\perp AB\)，得\(k_{AB}=1\)．\(AB\)过\(F(2,0)\)：\(y=x-2\)，即\(x-y-2=0\)，故选A．验算：\(P(-2,4)\)代入\(x-y+6=0\)：\(-2-4+6=0\)合．}
\fi
\end{ansblock}

\begin{ansblock}[2章-练-课时15-15]
% ans:2章-练-课时15-15
\ansitem{31}{(1) y²=4x；(2) 存在，此时 Q(9,6) 或 Q(4,−4)。}
\ifansbookdetail
\ansnote{详解}{(1)\(16=8p\)得\(p=2\)，\(C\)：\(y^{2}=4x\)．\par\noindent (2)\(F(1,0)\)．设\(N(t,2t+6)\)、\(Q(x_0,y_0)\)．四边形\(PQFN\)为平行四边形\(\Leftrightarrow\overrightarrow{PQ}=\overrightarrow{NF}\Leftrightarrow Q-P=F-N\)：\((4-t,-2t-2)=(x_0-1,y_0)\)，得\(x_0=5-t\)，\(y_0=-2t-2\)．\(Q\)在\(C\)上：\((-2t-2)^{2}=4(5-t)\)，即\(4t^{2}+8t+4=20-4t\)，\(t^{2}+3t-4=0\)，\(t=1\)或\(t=-4\)．\(t=1\)：\(N(1,8)\)，\(Q(4,-4)\)；\(t=-4\)：\(N(-4,-2)\)，\(Q(9,6)\)．\par\noindent 非退化核验（\(P\)、\(Q\)、\(F\)、\(N\)不共线，\(N\neq F\)）：\(t=1\)时\(\overrightarrow{PQ}=(0,-8)=\overrightarrow{NF}\)合；\(\overrightarrow{QF}=(-3,4)\)，\(\overrightarrow{PN}=\overrightarrow{P}-\overrightarrow{N}=(3,-4)\)，平行且相等合；四点不共线合，\(N(1,8)\neq F(1,0)\)合；\par\noindent \(t=-4\)时\(\overrightarrow{PQ}=(5,2)\)，\(\overrightarrow{NF}=\overrightarrow{F}-\overrightarrow{N}=(5,2)\)合；\(\overrightarrow{QF}=(-8,-6)\)，\(\overrightarrow{PN}=(8,6)\)合；斜率\(k_{PQ}=2/5\neq k_{QF}=3/4\)，四点不共线合，\(N\neq F\)合．两解均成立：存在\(N\)使\(PQFN\)为平行四边形，\(Q(9,6)\)或\(Q(4,-4)\)．验算：\(t^{2}+3t-4=(t-1)(t+4)\)合；\(Q(9,6)\)：\(36=4\times 9\)合；\(Q(4,-4)\)：\(16=16\)合．}
\fi
\end{ansblock}

\begin{ansblock}[2章-导-课时15-G1]
% ans:2章-导-课时15-G1
\ansitem{32}{B}
\ifansbookdetail
\ansnote{详解}{设\(M(x_0,y_0)\)．到对称轴的距离\(|y_0|=2\sqrt{2}\)，故\(y_0^{2}=8\)；到准线\(x=-p/2\)的距离\(x_0+p/2=3\)，得\(x_0=3-p/2\)．代入\(y_0^{2}=2px_0\)：\(8=2p(3-p/2)=6p-p^{2}\)，即\(p^{2}-6p+8=0\)，解得\(p=2\)或\(p=4\)．两值均有\(x_0=3-p/2\geqslant 0\)（\(p=2\)时\(x_0=2\)；\(p=4\)时\(x_0=1\)），皆合题意，故选B．验算：\(p=2\)时准线\(x=-1\)，\(M(2,\pm 2\sqrt{2})\)，\(x_0+1=3\)合；\(p=4\)时准线\(x=-2\)，\(M(1,\pm 2\sqrt{2})\)，\(x_0+2=3\)合．}
\fi
\end{ansblock}

\begin{ansblock}[2章-导-课时15-G2]
% ans:2章-导-课时15-G2
\ansitem{33}{B}
\ifansbookdetail
\ansnote{详解}{\(y=-(1/8)x^{2}\)即\(x^{2}=-8y\)，\(2p=8\)、\(p=4\)，开口向下，准线\(y=p/2=2\)，故选B．验算：焦点\((0,-2)\)、准线\(y=2\)，顶点\((0,0)\)到两者等距2合．}
\fi
\end{ansblock}

\begin{ansblock}[2章-导-课时15-G3]
% ans:2章-导-课时15-G3
\ansitem{34}{D}
\ifansbookdetail
\ansnote{详解}{两边同除以5：\(\sqrt{(x-1)^{2}+(y-2)^{2}}=|3x-4y+12|/5\)，而\(|3x-4y+12|/\sqrt{3^{2}+4^{2}}\)恰为\(M\)到直线\(3x-4y+12=0\)的距离．故动点\(M\)到定点\((1,2)\)的距离等于它到定直线\(3x-4y+12=0\)的距离；又\(3\times 1-4\times 2+12=7\neq 0\)，定点不在定直线上，由抛物线的定义，轨迹是以\((1,2)\)为焦点、\(3x-4y+12=0\)为准线的抛物线，故选D．说明：系数5的作用仅是把右端配成点到直线距离的标准形式（\(5=\sqrt{3^{2}+(-4)^{2}}\)），并非离心率记号．}
\fi
\end{ansblock}

\begin{ansblock}[2章-导-课时15-G4]
% ans:2章-导-课时15-G4
\ansitem{35}{y²=16x 或 x²=−8y}
\ifansbookdetail
\ansnote{详解}{直线\(x-2y-4=0\)与坐标轴的交点：令\(y=0\)得\(A(4,0)\)；令\(x=0\)得\(B(0,-2)\)．①焦点为\(A(4,0)\)：开口向右，设\(y^{2}=2px\)（\(p>0\)），\(p/2=4\)得\(p=8\)，方程\(y^{2}=16x\)；②焦点为\(B(0,-2)\)：开口向下，设\(x^{2}=-2py\)（\(p>0\)），\(p/2=2\)得\(p=4\)，方程\(x^{2}=-8y\)．综上，所求抛物线的方程为\(y^{2}=16x\)或\(x^{2}=-8y\)．验算：\(y^{2}=16x\)的焦点\((4,0)\)满足\(4-0-4=0\)，在直线上合；\(x^{2}=-8y\)的焦点\((0,-2)\)满足\(0+4-4=0\)，在直线上合．}
\fi
\end{ansblock}

\begin{ansblock}[2章-导-课时15-G5]
% ans:2章-导-课时15-G5
\ansitem{36}{y²=−8x；m=±2√6}
\ifansbookdetail
\ansnote{详解}{\(M(-3,m)\)的横坐标为负，抛物线开口向左，设\(y^{2}=-2px\)（\(p>0\)），则焦点\(F(-p/2,0)\)，准线\(x=p/2\)．由抛物线的定义，\(|MF|\)等于\(M\)到准线的距离：\(p/2+3=5\)，得\(p=4\)，故标准方程为\(y^{2}=-8x\)．将\(M(-3,m)\)代入：\(m^{2}=24\)，\(m=\pm 2\sqrt{6}\)．验算：焦点\(F(-2,0)\)：\(|MF|=\sqrt{(-3+2)^{2}+m^{2}}=\sqrt{1+24}=5\)合（两点距与定义两法互证）．}
\fi
\end{ansblock}

% ---- 尾块（\tailfill[30mm] 定高参——钉档 §三.3：无参在平衡栏呈「内容尾随框」，贴栏底须定高参；实测无参致末页平衡 Overfull\vbox 12.99pt，定高 30mm 三零） ----
\tailfill[30mm]

\end{multicols}
\end{document}
